\subsection{Deterministic covariance truncation}
\label{sec:appendix_deterministic_truncation}
\label{sec:appendix_necessity_sste}

Deterministic truncation exploits the spectral structure of the initial covariance, which the
Hermitian positive-definite matrix \(\mathsfbi{C}_{00}\in\mathbb{C}^{N\times N}\) supplies through
the eigendecomposition
\begin{equation}
  \mathsfbi{C}_{00}
  =
  \mathsfbi{U}\boldsymbol{\Sigma}\mathsfbi{U}^{*},
  \qquad
  \boldsymbol{\Sigma}
  =
  \operatorname{diag}
  \left(
    \sigma_1,\dots,\sigma_N
  \right),
  \qquad
  \sigma_1\geq\dots\geq\sigma_N>0,
  \label{eq:C00_eigendecomposition_sste_necessity}
\end{equation}
and which delivers the square-root factor \(\mathsfbi{L}\) in the form
\begin{equation}
  \mathsfbi{L}
  =
  \mathsfbi{U}\boldsymbol{\Sigma}^{1/2},
  \qquad
  \boldsymbol{\Sigma}^{1/2}
  =
  \operatorname{diag}
  \left(
    \sqrt{\sigma_1},\dots,\sqrt{\sigma_N}
  \right),
  \label{eq:C00_square_root_factor_sste_necessity}
\end{equation}
so that each eigenpair carries a single column of the factor into the dynamics.

DCT replaces \(\mathsfbi{C}_{00}\) by the rank-\(r\) truncation \(\mathsfbi{C}_{00}^{(r)}\) that
retains the leading \(r\) eigenpairs and propagates every retained column up to the target time
\(T\). \Cref{alg:deterministic_truncation} reports the resulting algorithm. The cost, in terms of
time-stepper calls, is
\begin{equation}
  \mathcal{C}_{\mathrm{DCT}}(T)=r,
  \label{eq:deterministic_fixed_time_cost}
\end{equation}
and, because we need to repeat this for all the \(N_t\) target times, the total cost is
\begin{equation}
  \mathcal{W}_{\mathrm{DCT}}=N_t\,r.
  \label{eq:deterministic_history_cost}
\end{equation}
Whenever the covariance has rank one, or an exact rank that a modest \(r\) already captures, DCT
returns \(g(T)\) exactly, and we prefer it to every alternative; whenever the spectrum extends
beyond the retained rank, DCT commits the error
\begin{equation}
  \varepsilon_r
  =
  \frac{1}{\tr\left\{\mathsfbi{W}\mathsfbi{C}_{00}\right\}}
  \tr
  \left\{
    \mathsfbi{T}
    \left(
      \mathsfbi{C}_{00}
      -
      \mathsfbi{C}_{00}^{(r)}
    \right)
  \right\}
  =
  \frac{1}{\tr\left\{\mathsfbi{W}\mathsfbi{C}_{00}\right\}}
  \sum_{j=r+1}^{N}
  \sigma_j\,
  \mathbf{u}_j^{*}\mathsfbi{T}\mathbf{u}_j.
  \label{eq:trace_error_1}
\end{equation}

\begin{algorithm}[H]
  \caption{Deterministic truncation (DCT) estimate of \(g(T)\)}
  \label{alg:deterministic_truncation}

  \begin{algorithmic}[1]
    \Require Initial covariance \(\mathsfbi{C}_{00}\), truncation rank \(r\),
             target time \(T\), operators \(\mathsfbi{A}(t)\) and \(\mathsfbi{B}(t)\),
             energy weight \(\mathsfbi{W}\)
    \Ensure Truncated growth functional \(g^{(r)}(T)\)

    \State Compute the eigenpairs \((\sigma_j,\mathbf{u}_j)\) of \(\mathsfbi{C}_{00}\)
      and order them so that \(\sigma_1\geq\dots\geq\sigma_N\)
      \Comment{Eigendecomposition}

    \State \(s\gets0\)
      \Comment{Trace accumulator}

    \For{\(j=1,\dots,r\)}
      \State \(\hat{\mathbf{q}}_j(0)\gets\sqrt{\sigma_j}\,\mathbf{u}_j\)
        \Comment{Retained column of \(\mathsfbi{L}\)}

      \State Solve
        \(\mathsfbi{B}(t)\,\partial_t\hat{\mathbf{q}}_j(t)
        =
        \mathsfbi{A}(t)\hat{\mathbf{q}}_j(t)\),
        \(0\leq t\leq T\)
        \Comment{Forward solve}

      \State \(s\gets s+\hat{\mathbf{q}}_j^{*}(T)\mathsfbi{W}\hat{\mathbf{q}}_j(T)\)
        \Comment{Energy of the propagated column}
    \EndFor

    \State \(g^{(r)}(T)\gets s/\tr\left\{\mathsfbi{C}_{00}\mathsfbi{W}\right\}\)
      \Comment{Normalize by the initial energy}

    \State \Return \(g^{(r)}(T)\)
  \end{algorithmic}
\end{algorithm}

\Cref{eq:trace_error_1} shows that two mechanisms govern the relative error. Choosing \(r\) such
that \(\sigma_r\simeq0\) makes the error small; a second mechanism nevertheless controls it, since
we can rearrange \eqref{eq:trace_error_1} into
\begin{equation}
  \varepsilon_r
  =
  \frac{1}{\tr\left\{\mathsfbi{W}\mathsfbi{C}_{00}\right\}}
  \sum_{j=r+1}^{N}
  \sigma_j
  \left\|\boldsymbol{\Phi}\mathbf{u}_j\right\|_{\mathsfbi{W}}^{2},
  \label{eq:trace_error_phi_tail}
\end{equation}
which requires \(\left\|\boldsymbol{\Phi}\mathbf{u}_j\right\|_{\mathsfbi{W}}^{2}\) to be small as
well.

This suggests that small neglected eigenvalues guarantee little on their own, since a discarded
eigendirection that aligns with a strongly amplified direction of the dynamics still contributes
substantially to the error. Consider in fact the simple case where
\begin{equation}
  \boldsymbol{\Phi}
  =
  \begin{pmatrix}
    1 & 0 \\
    0 & \sqrt{\kappa/\epsilon}
  \end{pmatrix},
  \qquad
  \mathsfbi{C}_{00}
  =
  \begin{pmatrix}
    1 & 0 \\
    0 & \epsilon
  \end{pmatrix},
  \qquad
  \mathsfbi{C}_{00}^{(r)}
  =
  \begin{pmatrix}
    1 & 0 \\
    0 & 0
  \end{pmatrix},
  \qquad
  \epsilon \ll 1 ,
  \label{eq:low_rank_failure_example}
\end{equation}
with \(\kappa\) the gain ratio and \(\epsilon\) the small covariance eigenvalue. Being
\(g^{(r)}\) the growth functional computed with \(\mathsfbi{C}_{00}^{(r)}\), the relative error is
\begin{equation}
  \frac{
    g(T)-g^{(r)}(T)
  }{
    g(T)
  }
  =
  \frac{\kappa}{\kappa+1}.
  \label{eq:low_rank_failure_relative_error}
\end{equation}
The discarded covariance eigenvalue \(\epsilon\) stays small while its contribution to the
propagated energy does not, precisely because the dynamics amplify the corresponding
eigendirection.

We therefore argue that a low-rank truncation of the initial covariance fails exactly when the
low-energy directions of \(\mathsfbi{C}_{00}\) intersect the high-gain directions of the evolution
operator, and that DCT delivers its substantial savings only when the covariance spectrum decays
rapidly and the dynamics stay benign. In this sense one does not know \emph{a priori} the value of
\(r\) that must be chosen, and reaching a reasonable estimate for covariances with a slowly
decaying spectrum may require \(r\sim N\), which leads to
\begin{equation}
  \mathcal{W}_{\mathrm{DCT}}=N\,N_t.
  \label{eq:deterministic_full_rank_cost}
\end{equation}

Summarizing, DCT carries two advantages, namely that it requires no adjoint problem, and that a
low-rank or rapidly decaying spectrum with \(\mathrm{rank}(\mathsfbi{C}_{00})\ll N\) keeps
\(r\sim O(1)\). It carries one disadvantage, namely that a correlation with a broad or full
spectrum, \(\mathrm{rank}(\mathsfbi{C}_{00})=N\), leads to large errors unless we choose the
directions of large gain carefully, and those directions are in principle not known \emph{a
priori}.
